\documentclass[10pt,a4paper]{article}

\usepackage[a4paper,margin=18mm]{geometry}
\usepackage[T1]{fontenc}
\usepackage[utf8]{inputenc}
\usepackage{charter}
\usepackage{amsmath,amssymb}
\usepackage{tabularx,multirow,array}
\usepackage{enumitem}
\usepackage[table]{xcolor}
\usepackage{titlesec}
\usepackage[hidelinks]{hyperref}

\setlist{nosep,leftmargin=14pt,itemsep=2pt}
\setlength{\parindent}{0pt}
\setlength{\parskip}{4pt}
\renewcommand{\arraystretch}{1.25}
\titleformat{\section}{\large\bfseries}{\thesection.}{0.5em}{}[\titlerule]
\titlespacing*{\section}{0pt}{10pt}{4pt}
\titleformat{\subsection}{\normalsize\bfseries}{}{0em}{}
\titlespacing*{\subsection}{0pt}{6pt}{2pt}

% Yellow-highlighted placeholder: remove once filled in.
\newcommand{\placeholder}[1]{\colorbox{yellow}{#1}}
\newcommand{\eps}{\varepsilon}

\begin{document}

%% ---------------------------------------------------------------- Overview
\thispagestyle{empty}
{\LARGE\bfseries Scaling Counterfactual Regret Minimization: a controlled study across poker game size and player count\par}
\vspace{4pt}
{\large CS F407 / U407 Artificial Intelligence --- Project Proposal\par}
\vspace{12pt}

\begin{tabularx}{\textwidth}{|>{\columncolor{gray!10}\bfseries}p{0.24\textwidth}|X|}
\hline
Project name & pokercfr \\ \hline
Team number / name & Team 29 --- \placeholder{<team name>} \\ \hline
Topic area & Agents \& Game Theory (imperfect-information games); Deep Learning used only in a gated stretch component (Deep CFR) \\ \hline
Track & \textbf{Research}: rigorous convergence / exploitability benchmarking. No live playable bot or UI. \\ \hline
GitHub repository & \url{https://github.com/Samanyu-dev/pokercfr} (creator: Samanyu-dev) \\ \hline
Team members &
Samanyu Reddy Allipuram (2023A2PS0910H) \newline
Himnish Lalchandani (2023AAPS1131H) \newline
Aarush Kanipakam (2023A7PS1160H) \newline
Sai Aashish Juryala (2023A7PS0146H) \newline
Siddhant Narlawar (2023A7PS0049H) \newline
Shashwat Kumar (2023A7PS1235H) \newline
Akhil Thindi (2023AAPS0172H) \newline
Vikranth \placeholder{<surname> (<ID>)} \\ \hline
Short description &
We do not propose a new algorithm. We build one parameterized family of poker games, from Kuhn Poker up through Leduc Hold'em and multiplayer Leduc, where exact best responses are computable at every size. On it we measure how the cost of reaching an $\eps$-Nash equilibrium grows for CFR, CFR+, Discounted CFR and Monte Carlo CFR as we change \emph{one} factor at a time: deck size, cards dealt, betting structure, or number of players. We also measure how far CFR is from equilibrium in 3--6 player games, where its convergence guarantee no longer holds. ``Solved'' is defined quantitatively as exact exploitability (2 players) or NashConv ($N$ players) below a pre-registered threshold. \\ \hline
\end{tabularx}

\newpage
\setcounter{page}{1}

%% ---------------------------------------------------------------- 1
\section{Problem statement \& motivation}

Counterfactual Regret Minimization (CFR)~[1] and its descendants~[2--5] power every superhuman poker agent to date~[7--9]. Each of these results, however, is reported on one game at a time: Kuhn, Leduc, or Hold'em, with deck size, hand size, betting structure and player count all changing together. A practitioner therefore cannot answer a basic scaling question: \emph{which} property of a game makes it expensive for CFR, and does the relative ranking of CFR variants survive as games grow? A second open gap is multiplayer play. CFR's Nash convergence guarantee holds only for two-player zero-sum games, yet Pluribus~[9] succeeded at six players. How far CFR actually lands from equilibrium as player count grows is rarely measured, because exact best responses are intractable in real poker.

\textbf{Novelty} (responding to the ``Zinkevich 2007 is exactly this'' feedback). We reproduce CFR only as a validated baseline. The contribution is a \textbf{controlled scaling study} built on a game family we implement ourselves, in which deck size, cards dealt and player count are \emph{independent} parameters. Standard multiplayer Leduc implementations grow the deck with the number of players (OpenSpiel~[11] uses $2(N+1)$ cards), which confounds the two effects. Because every game in our family is small enough for exact best-response computation, every claim is checked against ground truth rather than self-play win rates.

\subsection{Research hypotheses}
\begin{itemize}
  \item \textbf{H1 (primary, pre-registered).} At equal wall-clock budgets, CFR+'s average strategy has lower exact exploitability than vanilla CFR's on two-player Leduc Hold'em, across pre-specified raise caps.
  \item \textbf{H2 (scaling).} Iterations and wall-clock time to reach a fixed $\eps$ grow at measurably different rates along the deck-size, cards-dealt and betting-depth axes. The CFR-family ranking found in H1 is either preserved across sizes or breaks at an identifiable point.
  \item \textbf{H3 (players).} For $N \ge 3$ players, CFR's NashConv decreases but plateaus above zero, and the plateau rises with $N$ while the deck is held fixed.
\end{itemize}

%% ---------------------------------------------------------------- 2
\section{Connection to course concepts}
\begin{itemize}
  \item \textbf{Extensive-form games with imperfect information}: game trees with chance nodes, information sets (states a player cannot distinguish), behavioural mixed strategies.
  \item \textbf{Nash equilibrium, $\eps$-Nash equilibrium and best response}. In two-player zero-sum games the minimax theorem makes equilibrium value unique. Kuhn Poker's closed-form value is $-1/18$ for the first player~[10].
  \item \textbf{Regret minimization / no-regret learning}: regret matching at each information set; average-strategy convergence to equilibrium in two-player zero-sum self-play~[1]. Variants include regret-matching+ (CFR+)~[2] and discounted regrets (DCFR)~[4].
  \item \textbf{Game-tree search}: depth-first recursive traversal computing expected values over chance and opponent nodes (expectimax-style). Best response is computed by a max over the responder's information sets.
  \item \textbf{Monte Carlo sampling}: outcome/chance-sampled MCCFR~[3] for games too large to traverse fully each iteration.
  \item \textbf{Deep learning (stretch only)}: Deep CFR~[5] replaces regret tables with an MLP trained by regression on a reservoir-sampled buffer.
\end{itemize}

%% ---------------------------------------------------------------- 3
\section{Proposed approach \& methodology}

\subsection{Environment}
All games are implemented in our own Python engine; no poker data is needed. One configuration, \texttt{GameConfig(ranks, suits, cards\_dealt, players)}, generates the whole Kuhn family. The betting is Kuhn-style $N$-player betting: ante 1, check or bet 1, and after a bet every other player calls or folds once. With more than one card dealt, hands are ranked by pairs, two pair, trips and quads, then by high cards. Straights and flushes are added only if we reach 5-card hands. A separate Leduc engine adds a public card, two betting rounds and a configurable raise cap. The ladder engine and an $N$-player vanilla CFR are already in the repository and unit-tested: CFR reproduces Kuhn's $-1/18$ value and its known equilibrium family.

\subsection{The game ladder: what we will solve, smallest first}
We grow the game one factor at a time. Each step changes a single parameter of the previous one, so its effect on solving cost can be isolated (H2, H3). Sizes below are computed by our engine. A ``deal'' is one way of distributing private cards; full-traversal CFR visits every deal on every iteration.

\begin{center}\small
\begin{tabularx}{\textwidth}{|l|X|l|l|l|}
\hline
\rowcolor{gray!15}\textbf{Step} & \textbf{Game} & \textbf{Factor changed} & \textbf{Info sets} & \textbf{Deals} \\ \hline
1 & Kuhn: J Q K, 1 card each, 2 players & --- (validation) & 12 & 6 \\ \hline
2 & J Q K A, 1 card, 2 players & deck: +1 rank & 16 & 12 \\ \hline
3 & J Q K A, 2 cards, 2 players & cards dealt $1\to2$ & 24 & 6 \\ \hline
4 & J Q K A $\times$ 2 suits (8 cards), 1 card & deck: +1 suit & 32 & 56 \\ \hline
5 & 8 cards, $2\to3\to4$ cards dealt & cards dealt & $112\to280$ & $420\to70$ \\ \hline
6 & 8 cards, 1 card, $3\to4$ players & players $2\to3\to4$ & $96\to256$ & $336\to1{,}680$ \\ \hline
7 & 8 cards, 2 cards, 3 players; 12 cards, 2 cards, 4 players & players $\times$ cards & $336\to2{,}112$ & $2{,}520\to1.2$\,M \\ \hline
8 & 24 cards (6 ranks $\times$ 4 suits), 2 cards, 2 players & deck size & 1,104 & 63,756 \\ \hline
9 & 24 cards, 3 cards, 3 players $\to$ 4 cards, 4 players & all three & 24\,K $\to$ 340\,K & $2.2{\times}10^{9} \to 4.6{\times}10^{13}$ \\ \hline
10 & Full 52-card deck, 2 cards, 2 players & deck size & 5,304 & 1.6\,M \\ \hline
11 & Leduc Hold'em, 2 players; raise caps varied & public card, 2 rounds & \multicolumn{2}{l|}{main testbed for H1} \\ \hline
12 & Leduc Hold'em, $3\to6$ players, deck held fixed & players & \multicolumn{2}{l|}{H3} \\ \hline
13 & \textit{(Stretch)} Heads-up \emph{limit} Hold'em with card abstraction; Deep CFR on Leduc & scale & \multicolumn{2}{l|}{gated, see below} \\ \hline
\end{tabularx}
\end{center}

Two observations drive our solver choice. First, with one betting round the number of information sets stays small even with a full deck; cost comes from the number of deals. Steps 1--8 are solved by exact full-traversal CFR. From step 9 on (and step 10 at scale), deals grow far faster than information sets, so we switch to chance-sampled MCCFR~[3] and to a vectorized traversal that evaluates all private hands at once. We report the step at which each method stops being feasible on our hardware. This ``how far can each method go'' result is itself one of our findings.

\textbf{Scope, stated explicitly.} Steps 1--12 are the semester target. Step 13 starts only if steps 1--12 have complete results by week 6. Full (no-limit, multi-player) Texas Hold'em is \textbf{out of scope}. Even heads-up limit Hold'em has ${\sim}3{\times}10^{14}$ information sets and took ${\sim}900$ core-years to solve exactly~[7], so step 13 would only ever use a coarse abstraction and no hypothesis depends on it.

\subsection{Algorithms}
Vanilla CFR~[1], CFR+~[2], Discounted CFR~[4] and outcome-sampling MCCFR~[3], all sharing one traversal and evaluation harness. A study cell is an (algorithm $\times$ game $\times$ seed) run, and every cell reports exploitability against both iterations and wall-clock time.

\subsection{Evaluation: what ``solved'' means}
\begin{itemize}
  \item \textbf{Primary metric, two players:} exact exploitability
  \[ e(\sigma) = \tfrac{1}{2}\Big[\max_{\sigma_1'} u_1(\sigma_1', \sigma_2) + \max_{\sigma_2'} u_2(\sigma_1, \sigma_2')\Big], \]
  computed by a full best-response traversal and reported in chips/hand. A game counts as \textbf{solved to $\eps$} when $e(\sigma) \le \eps$, with $\eps$ fixed per game \emph{before} final runs.
  \item \textbf{Primary metric, $N$ players:} $\mathrm{NashConv}(\sigma) = \sum_i \big[\max_{\sigma_i'} u_i(\sigma_i', \sigma_{-i}) - u_i(\sigma)\big]$, i.e.\ the total gain available to unilateral deviators.
  \item \textbf{Ground-truth checks:} Kuhn's closed-form equilibrium family~[10]. On Kuhn and two-player Leduc, our exploitability numbers are cross-checked against OpenSpiel's independent implementation~[11].
  \item \textbf{Secondary:} head-to-head results against the opponent pool below, using fixed seeds with 95\% confidence intervals. These are reported as supporting evidence only, never as a measure of equilibrium quality.
  \item \textbf{Rigor:} $\ge 5$ seeds per cell, pre-registered compute budgets and raise caps, and a fresh-clone reproducibility check before submission.
\end{itemize}

\subsection{Who our agent plays against}
For every step, the trained agent is evaluated against a fixed pool of opponents, from weakest to strongest:
\begin{enumerate}
  \item \textbf{Fixed-policy baselines}: uniform random, always-call, always-bet/raise, and a rule-based agent that bets only above a hand-strength threshold.
  \item \textbf{Cross-play between our solvers}: CFR vs CFR+ vs DCFR vs MCCFR agents trained under the same compute budget, so differences in exploitability can be compared with differences in actual play.
  \item \textbf{External reference}: OpenSpiel's CFR agents on Kuhn and Leduc~[11].
  \item \textbf{The best-response opponent}: the strongest possible counter-strategy to our agent, computed exactly. This is the worst case, and its winnings are the exploitability above.
\end{enumerate}
Matches use fixed seeds with seat rotation (each deal replayed with positions swapped) to cut card-luck variance. Results are reported in chips/hand with 95\% confidence intervals. We also track how the agent's bluffing and calling frequencies change at each step, compared with Kuhn's known equilibrium values.

\subsection{Compute}
Steps 1--12 run on team laptops (Apple M-series, 24\,GB RAM). Measured: one vanilla CFR iteration takes 0.3\,ms on Kuhn and 0.5\,s on the 8-card, 4-player game. Independent experiment cells (step $\times$ algorithm $\times$ seed) are spread one per machine, since CFR tables need shared memory and are not split across machines. Kaggle notebooks (GPU) are used for Deep CFR, with checkpointing to survive the 12-hour session limit. AWS SageMaker is used only for runs that exceed laptop memory (large-deck MCCFR, the step-13 abstraction).

%% ---------------------------------------------------------------- 4
\section{Ethics}
Stronger poker AI is dual-use. It could be deployed for real-money online play, where bots violate platform rules and harm human players, and better agents are harder to detect or exploit. We bound the project as follows:
\begin{itemize}
  \item \textbf{Toy games only.} Kuhn, its variants and Leduc have no real-money market. The step-13 stretch targets heads-up \emph{limit} Hold'em, already essentially solved and publicly documented~[7], so the project creates no new real-world capability.
  \item \textbf{No deployment surface.} No interface to any poker site, no screen-reading or client automation, no live-play bot, and no work on human-like behaviour or detection evasion. The Research track was chosen partly for this reason.
  \item \textbf{No human subjects or personal data.} All evaluation is agent-vs-agent in simulation.
  \item \textbf{Transparency.} The repository README carries an intended-use and ethics statement, and the final report discloses all AI-assistant use and how outputs were verified, as the course requires.
\end{itemize}

%% ---------------------------------------------------------------- 5
\section{Current work in the area}
CFR~[1] introduced counterfactual regret and proved $O(1/\sqrt{T})$ convergence in two-player zero-sum games. CFR+~[2] and DCFR~[4] improve practical speed; MCCFR~[3] trades per-iteration cost for variance; Deep CFR~[5] removes tabular storage. At scale, these techniques solved heads-up limit Hold'em~[7] and beat professionals at heads-up no-limit (DeepStack~[12], Libratus~[8]) and six-player no-limit (Pluribus~[9]). For multiplayer settings, Abou Risk \& Szafron~[6] applied CFR to 3-player poker without equilibrium guarantees. Leduc Hold'em~[13] is the standard mid-size benchmark, and OpenSpiel~[11] provides reference implementations that we use for cross-validation, not as our engine. What we add is a controlled, one-factor-at-a-time scaling comparison with exact evaluation, not a new solver.

%% ---------------------------------------------------------------- 6
\section{Feasibility: timeline \& roles}
Final submission dates are TBA, so the timeline assumes ${\sim}9$ weeks from 5 Oct.

\begin{center}\small
\begin{tabularx}{\textwidth}{|l|X|}
\hline
\rowcolor{gray!15}\textbf{Week} & \textbf{Milestone} \\ \hline
1 (5--11 Oct) & Merge ladder engine + vanilla CFR; exploitability / NashConv calculator; step 1 (Kuhn) validated against closed form and OpenSpiel \\ \hline
2 (12--18 Oct) & CFR+; steps 2--5 solved; Leduc environment + tests; baseline bots \\ \hline
3 (19--25 Oct) & DCFR, MCCFR; steps 6--8; seeding/config system; benchmarking harness and plotting pipeline \\ \hline
4 (26 Oct--1 Nov) & Pre-register $\eps$, budgets and raise caps; full sweep of steps 1--8 (H2); two-player Leduc runs (H1) \\ \hline
5 (2--8 Nov) & Vectorized traversal + MCCFR on steps 9--10; multiplayer Leduc, $N = 3$--$6$ at a fixed deck (H3); mid-project check-in \\ \hline
6 (9--15 Nov) & Opponent-pool tournament across all steps; \textbf{stretch gate decision} \\ \hline
7 (16--22 Nov) & Multi-seed runs + statistics; stretch work if the gate passed (Deep CFR / abstracted limit Hold'em) \\ \hline
8 (23--29 Nov) & Ablations; fresh-clone reproducibility check; final plots \\ \hline
9 (30 Nov--6 Dec) & Report, AI-usage disclosure, demo video, viva preparation \\ \hline
\end{tabularx}

\vspace{6pt}
\begin{tabularx}{\textwidth}{|l|l|X|}
\hline
\rowcolor{gray!15}\textbf{Layer} & \textbf{Member} & \textbf{Owns} \\ \hline
\multirow{2}{*}{CFR engine} & Himnish Lalchandani & Vanilla CFR, CFR+, DCFR; code review lead \\ \cline{2-3}
 & Vikranth & MCCFR; Deep CFR (stretch) \\ \hline
\multirow{3}{*}{\shortstack[l]{Game state \&\\abstraction}} & Samanyu Reddy Allipuram & Leduc engine, raise caps, multiplayer Leduc \\ \cline{2-3}
 & Sai Aashish Juryala & Kuhn-family ladder engine (steps 1--10), environment tests, seeding/reproducibility \\ \cline{2-3}
 & Siddhant Narlawar & Vectorized traversal for large decks, compute runs on Kaggle/AWS, card abstraction (stretch) \\ \hline
\multirow{3}{*}{\shortstack[l]{Evaluation \&\\reporting}} & Aarush Kanipakam & Exploitability / NashConv, experiment design, results matrix, report lead \\ \cline{2-3}
 & Shashwat Kumar & Baseline bots, opponent-pool tournament harness \\ \cline{2-3}
 & Akhil Thindi & Statistics (multi-seed, confidence intervals), plots and strategy visualizations, fresh-clone reproducibility check \\ \hline
\end{tabularx}
\end{center}

Every change goes through a GitHub issue and a reviewed pull request, so the commit history records each member's contribution across the semester.

%% ---------------------------------------------------------------- References
\section*{References}
\sloppy
\begin{enumerate}[label={[\arabic*]},leftmargin=22pt,itemsep=1pt]\small
  \item M.~Zinkevich, M.~Johanson, M.~Bowling, C.~Piccione. Regret Minimization in Games with Incomplete Information. NeurIPS 2007.
  \item O.~Tammelin. Solving Large Imperfect Information Games Using CFR+. arXiv:1407.5042, 2014.
  \item M.~Lanctot, K.~Waugh, M.~Zinkevich, M.~Bowling. Monte Carlo Sampling for Regret Minimization in Extensive Games. NeurIPS 2009.
  \item N.~Brown, T.~Sandholm. Solving Imperfect-Information Games via Discounted Regret Minimization. AAAI 2019. arXiv:1809.04040.
  \item N.~Brown, A.~Lerer, S.~Gross, T.~Sandholm. Deep Counterfactual Regret Minimization. ICML 2019. arXiv:1811.00164.
  \item N.~Abou Risk, D.~Szafron. Using Counterfactual Regret Minimization to Create Competitive Multiplayer Poker Agents. AAMAS 2010.
  \item M.~Bowling, N.~Burch, M.~Johanson, O.~Tammelin. Heads-up Limit Hold'em Poker is Solved. \emph{Science} 347(6218), 2015. doi:10.1126/science.1259433.
  \item N.~Brown, T.~Sandholm. Superhuman AI for Heads-up No-limit Poker: Libratus Beats Top Professionals. \emph{Science} 359(6374), 2018. doi:10.1126/science.aao1733.
  \item N.~Brown, T.~Sandholm. Superhuman AI for Multiplayer Poker. \emph{Science} 365(6456), 2019. doi:10.1126/science.aay2400.
  \item H.~W.~Kuhn. A Simplified Two-Person Poker. \emph{Contributions to the Theory of Games I}, Annals of Mathematics Studies 24, 1950.
  \item M.~Lanctot et al. OpenSpiel: A Framework for Reinforcement Learning in Games. arXiv:1908.09453, 2019.
  \item M.~Morav\v{c}\'{i}k et al. DeepStack: Expert-Level Artificial Intelligence in Heads-up No-limit Poker. \emph{Science} 356(6337), 2017. doi:10.1126/science.aam6960.
  \item F.~Southey et al. Bayes' Bluff: Opponent Modelling in Poker. UAI 2005.
\end{enumerate}

\end{document}
